\documentclass[10pt,a4paper]{amsart}
\usepackage[margin=19mm]{geometry}
\usepackage{amsmath,amssymb,mathtools}
\usepackage[scr=rsfs,cal=euler]{mathalfa}
\usepackage{xfrac}
\usepackage[unicode,hidelinks,bookmarks=false]{hyperref}
\newcommand{\ie}{i.e.\ }
\newcommand{\cf}{cf.\ }
\newcommand{\resp}{respectively\ }
\newcommand{\Subsection}{\subsection}
\providecommand{\nobreakdash}{\nobreak\hbox{-}}
\setlength{\emergencystretch}{2em}
\multlinegap0pt
\usepackage[shortlabels]{enumitem}
\usepackage{amssymb}
\usepackage{tcolorbox}
\newtheorem{lemma}{Lemma}
\newtheorem{theorem}{Theorem}
\def\F{{\mathbb F}}
\DeclareMathOperator{\rad}{rad}
\title{Nonlinear maps preserving the polynomial}
\author{Andrey Yurkov}
\date{Source: arXiv:2604.23690v1 (CC BY 4.0)}
\begin{document}
\maketitle

\noindent\emph{Source and licence.} Nonlinear maps preserving the polynomial. Version arXiv:2604.23690v1 (CC BY 4.0), distributed by arXiv under the Creative Commons Attribution 4.0 International licence, http://creativecommons.org/licenses/by/4.0/, as recorded in the arXiv licence metadata for this exact version and shown on the abstract page https://arxiv.org/abs/2604.23690v1. Full licence notice: \texttt{licenses/arxiv-2604.23690-CC-BY-4.0.txt}.\par\smallskip

\noindent\textit{Editorial scope.} Counted unit: the article's theorem thm:homogenphipsi (source0.tex lines 889--897). The article's complete original proof of it is delivered with it (source0.tex lines 898--1062, one \texttt{proof} environment of 165 lines). No mathematics is written, repaired or re-derived: the statement and the proof are the article's own lines, in the article's own notation, and no retained line is substituted or re-flowed.  Retained uncounted as the source-local context the proof needs: lines 247--256; lines 314--331; lines 712--717; lines 722--731.  Nothing was omitted from any retained span, so the package carries no omission record.  No retained line contains a citation, so the wrapper carries no bibliography and no bibliography entry.  No retained line contains a figure environment, an image-inclusion command or drawing code, so no image is delivered and the package declares no asset dependency.  Editorial formatting only, in addition: an AMS \texttt{amsart} preamble that restates the article's own declarations and macros used by the retained lines, namely the environments \texttt{lemma}, \texttt{theorem} on separate counters, together with the standard packages those lines need.  The retained environments print in this document as Lemma 1, Theorem 1 in order of appearance, whereas the article prints the same items with the numbering of its own full document.  The complete native source of the article as distributed in the arXiv e-print for this exact version is bundled unchanged as \texttt{sources/199181/source0.tex}.
\begin{lemma}\label{lem:StrangeCondImplyRadP}
Assume that $|\F| > 2$, $V$ is vector space over $\F$ and $f\colon V \to \F$ be a function. Let $\mathbf x \in V$. If  
\begin{equation}\label{lem:StrangeCondImplyRadP:eqq1}
f(\mathbf a + \lambda \mathbf x) = f(\mathbf a + \mathbf x)\quad\mbox{for all}\;\;\mathbf a \in V,\, 0 \neq \lambda \in \F, 
\end{equation}
then $\mathbf x \in \rad (f)$.
%\begin{equation}\label{lem:StrangeCondImplyRadP:eqq2}
%f(\mathbf a + \lambda \mathbf x) = f(\mathbf a)\quad\mbox{for all}\;\;\mathbf a \in \F^n,\, 0 \neq \lambda \in \F, 
%\end{equation}
\end{lemma}

\begin{lemma}\label{lem:TwoCondLift}
Let $V$ be a vector space over $\F$ and $f\colon V \to \F$ be a function, $\phi, \psi \colon V \to V$ be maps such that
\begin{equation}\label{lem:TwoCondLift:eqq1}
f(\mathbf x + \lambda \mathbf y) = f(\phi(\mathbf x) + \lambda (\mathbf y))\quad\mbox{for all}\;\;\mathbf x, \mathbf y \in V, \lambda \in \F,
\end{equation}
$\psi_{\rad(f)}\colon V/\rad(f) \to V/\rad(f)$ be a map satisfying
\begin{equation}\label{lem:TwoCondLift:eqq2}
\psi_{\rad(f)} \circ \pi_{\rad(f)} = \pi_{\rad(f)} \circ \psi. 
\end{equation}
Then
\begin{equation*}\label{lem:TwoCondLift:eqq3}
f_{\rad(f)}(\pi_{\rad(f)}(\mathbf x) + \lambda \mathbf y) = f_{\rad(f)}(\pi_{\rad(f)}(\phi(\mathbf x)) + \lambda \psi_{\rad(f)}(\mathbf y))
\end{equation*}
for all $\mathbf x \in V, \mathbf y \in V/\rad(f), \lambda \in \F.$
%\begin{equation}
%f(\mathbf a + \lambda \mathbf x) = f(\mathbf a)\quad\mbox{for all}\;\;\mathbf a \in \F^n,\, 0 \neq \lambda \in \F, 
%\end{equation}
\end{lemma}

\begin{lemma}\label{lem:condimpliesphi}Let $P \in \F[x_1,\ldots, x_n]$, and let $\phi, \psi \colon \F^n \to \F^n$ be two maps such that
\begin{equation}\label{lem:condimpliesphi:eqq1}
P(\mathbf{x} + \lambda \mathbf{y}) = P(\phi(\mathbf{x}) + \lambda \psi(\mathbf{y}))\;\;\mbox{for all}\;\;\mathbf{x},\mathbf{y} \in \mathbb \F^n,\; \lambda \in \F.
\end{equation}
Then $P(\phi(\mathbf x)) = P (\mathbf x)$ for all $\mathbf x \in \F^n.$
\end{lemma}

\begin{lemma}\label{lem:PPhiPsi}Let $P \in \F[x_1,\ldots, x_n]$ be a polynomial such that $\deg (P) < |\F|$ and $\dim (\mathcal L_P) + \dim (\rad(P)) = n$.  If  $\phi, \psi \colon \F^n \to \F^n$ are two maps such that
\begin{equation}\label{lem:PPhiPsi:eqq1}
P(\mathbf{x} + \lambda \mathbf{y}) = P(\phi(\mathbf{x}) + \lambda \psi(\mathbf{y}))\;\;\mbox{for all}\;\;\mathbf{x},\mathbf{y} \in \mathbb \F^n,\; \lambda \in \F,
\end{equation}
then there exists a unique bijective linear map $\psi_{\rad(P)}$ on $\F^n / \rad(P)$ such that 
\begin{equation}\label{lem:PPhiPsi:eq}
\psi_{\rad(P)} \circ \pi_{\rad(P)} = \pi_{\rad(P)} \circ \psi,
\end{equation}
 where $\pi_{\rad(P)}$ is a canonical projection on $\F^n / \rad(P)$.
\end{lemma} 

\begin{theorem}\label{thm:homogenphipsi}Let $P \in \F[x_1,\ldots, x_n]$ be a homogeneous polynomial such that $\deg (P) < |\F|$ and $\dim (\mathcal L_P) + \dim (\rad(P)) = n$.  If  $\phi, \psi \colon \F^n \to \F^n$ are such that
\begin{equation}\label{thm:homogenphipsi:eqq1}
P(\mathbf{x} + \lambda \mathbf{y}) = P(\phi(\mathbf{x}) + \lambda \psi(\mathbf{y}))\;\;\mbox{for all}\;\;\mathbf{x},\mathbf{y} \in \mathbb \F^n,\; \lambda \in \F,
\end{equation}
then there exists a unique linear map $T_{\rad} \colon \F^n/\rad(P) \to \F^n/\rad(P)$ preserving $P_{\rad(P)}$ such that
\begin{equation*}\label{thm:homogenphipsi:eq}
\pi_{\rad(P)} \circ \phi = \pi_{\rad(P)} \circ \psi = T_{\rad(P)} \circ \pi_{\rad(P)}. 
\end{equation*}
\end{theorem} 

\begin{proof}Let $P \in \F[x_1,\ldots, x_n]$ be polynomial and $\phi, \psi\colon \F^n \to \F^n$ be maps satisfying all the conditions of the lemma.  Let $\psi_{\rad}\colon \F^n / \rad(P) \to \F^n / \rad(P)$ be a bijective linear map satisfying the equality~\eqref{lem:PPhiPsi:eq}. Such a linear map exists by Lemma~\eqref{lem:PPhiPsi}. The following five claims will be proven sequently:
\begin{enumerate}[label=(\roman*), ref=(\roman*)]
\item\label{thm:homogenphipsi:part1} If $\deg(P) > 1$, then $\phi(\mathbf 0) \in \rad(P)$.
\item\label{thm:homogenphipsi:part2} $\psi$ preserves $P$.
\item\label{thm:homogenphipsi:part3} $\psi_{\rad(P)}$ preserves $P_{\rad}$.
\item\label{thm:homogenphipsi:part4} $\pi_{\rad(P)} \circ \phi = \pi_{\rad(P)} \circ \psi$.
\item\label{thm:homogenphipsi:part5} $\psi_{\rad(P)}$ is a unique linear map satisfying the conditions of the lemma.
\end{enumerate}
The claim~\ref{thm:homogenphipsi:part5} clearly implies the statement of the lemma.
\paragraph{Claim~\ref{thm:homogenphipsi:part1}}%. If $\deg(P) > 1$, then $\phi(\mathbf 0) \in \rad(P)$}
Assume in that  $d_P = \deg(P) > 1$. Let us ensure that $\phi(\mathbf 0) \in \rad(P)$ using Lemma~\ref{lem:StrangeCondImplyRadP}. For this, let us show that
\begin{equation}\label{thm:homogenphipsi:eq77}
P(\mathbf a + \lambda \phi (\mathbf 0)) = P(\mathbf a + \phi (\mathbf 0))\;\;\mbox{for all}\;\;\mathbf a \in \F^n\;\;\mbox{and}\;\;0\neq \lambda \in \F.
\end{equation}

Let $\mathbf a \in \F^n$ and $0 \neq \lambda \in \F$. Since we assume that $P$ is homogeneous,
\begin{equation}\label{thm:homogenphipsi:eq5}
P\left(\mathbf a + \lambda \phi (\mathbf 0)\right) = \lambda^{d_P} P\left(\frac{1}{\lambda}\mathbf a + \phi (\mathbf 0)\right).
\end{equation}
The map $\psi_{\rad(P)}$ is bijective. Consequently, there exists $\mathbf b \in \F^n$ such that 
\begin{equation}\label{thm:homogenphipsi:eq1}
\psi_{\rad(P)}\left(\pi_{\rad(P)}(\mathbf b)\right) = \pi_{\rad(P)}(\mathbf a).
\end{equation}
Hence,
\begin{equation*}
\pi_{\rad(P)} \left(\psi(\mathbf b)\right) \overset{\eqref{lem:PPhiPsi:eq}}{=\joinrel=} \psi_{\rad(P)}\left(\pi_{\rad(P)}(\mathbf b)\right)  \overset{\eqref{thm:homogenphipsi:eq1}}{=\joinrel=} \pi_{\rad(P)}(\mathbf a),
\end{equation*}
or, equivalently
\begin{equation}\label{thm:homogenphipsi:eq2}
\left(\psi (\mathbf b) - \mathbf a\right) \in \rad(P). 
\end{equation}
Therefore,
\begin{multline}\label{thm:homogenphipsi:eq6}
P(\frac{1}{\lambda}\mathbf a + \phi (\mathbf 0))\\ \overset{\mbox{\scriptsize The definition of}\;\rad(P)\;\mbox{\scriptsize and~\eqref{thm:homogenphipsi:eq2}}}{=\joinrel=\joinrel=\joinrel=\joinrel=\joinrel=\joinrel=\joinrel=\joinrel=\joinrel=\joinrel=\joinrel=\joinrel=\joinrel=\joinrel=\joinrel=\joinrel=\joinrel=\joinrel=} P\left(\frac{1}{\lambda} \left(\psi (\mathbf b) - \mathbf a\right) + \frac{1}{\lambda}\mathbf a + \phi (\mathbf 0)\right)\\ = P\left(\frac{1}{\lambda}\psi (\mathbf b) + \phi (\mathbf 0)\right)
\end{multline}
Since $\phi$ and $\psi$ satisfy the condition~\eqref{thm:homogenphipsi:eqq1}, then
\begin{equation}\label{thm:homogenphipsi:eq3}
P\left(\frac{1}{\lambda}\psi (\mathbf b) + \phi (\mathbf 0)\right) = P\left(\frac{1}{\lambda}\mathbf b + \mathbf 0\right) = \frac{1}{\lambda^{d_P}}P\left(\mathbf b\right)
\end{equation}
and
\begin{multline}\label{thm:homogenphipsi:eq4}
P\left(\mathbf b\right) = P\left(\mathbf 0 + \mathbf b\right) = P\left(\phi(\mathbf 0) + \psi(\mathbf b)\right)\\ \overset{\mbox{\scriptsize The definition of}\;\rad(P)\;\mbox{\scriptsize and~\eqref{thm:homogenphipsi:eq2}}}{=\joinrel=\joinrel=\joinrel=\joinrel=\joinrel=\joinrel=\joinrel=\joinrel=\joinrel=\joinrel=\joinrel=\joinrel=\joinrel=\joinrel=\joinrel=\joinrel=\joinrel=\joinrel=}
P\left(\phi(\mathbf 0) + \psi(\mathbf b) + \left(\mathbf a - \psi(\mathbf b)\right)\right)\\
=  P\left(\phi(\mathbf 0) + \mathbf a\right).
\end{multline}


Thus,
\begin{multline*}
P\left(\mathbf a + \lambda \phi (\mathbf 0)\right) \overset{\eqref{thm:homogenphipsi:eq5}}{=\joinrel=}  \lambda^{d_P} P\left(\frac{1}{\lambda}\mathbf a + \phi (\mathbf 0)\right)\overset{\eqref{thm:homogenphipsi:eq6}}{=\joinrel=}\\
\lambda^{d_P} P\left(\frac{1}{\lambda}\psi (\mathbf b) + \phi (\mathbf 0)\right) \overset{\eqref{thm:homogenphipsi:eq3}}{=\joinrel=} \lambda^{d_P} \cdot \frac{1}{\lambda^{d_P}}P\left(\mathbf b\right) \overset{\eqref{thm:homogenphipsi:eq4}}{=\joinrel=} P\left(\phi(\mathbf 0) + \mathbf a\right),
\end{multline*}
which holds for all $\mathbf a \in \F^n$ and $0 \neq \lambda \in \F$. From this we conclude that~\eqref{thm:homogenphipsi:eq77} holds for $\phi(\mathbf 0)$. Since $\deg(P) > 1$, then $|\F| > \deg(P)$ implies that $|\F| > 2$. Hence, the conditions Lemma~\ref{lem:StrangeCondImplyRadP} are satisfied for $\F$ and $\phi(\mathbf 0)$  and consequently $\phi(\mathbf 0) \in \rad(P).$

\paragraph{Claim~\ref{thm:homogenphipsi:part2}}%. $\psi$ preserves $P$}
Let us consider three cases: (I) $\deg(P) < 1$; (II) $\deg(P) = 1$ and (III) $\deg(P) > 1$.

\paragraph{Case (I): $\deg(P) < 1$} Then $P$ is a constant function. Therefore, $\phi$ tautologically preserves $P$. 

\paragraph{Case (II): $\deg(P) = 1$} Then $P$ is a linear function. Therefore, the condition~\eqref{thm:homogenphipsi:eqq1} can be rewritten as follows
\begin{multline}\label{thm:homogenphipsi:eq7}
P(\mathbf x) + \lambda P(\mathbf y) = P(\mathbf x + \lambda \mathbf y) = P(\phi(\mathbf x) + \lambda \psi(\mathbf y))\\ = P(\phi(\mathbf x)) + \lambda P(\psi(\mathbf y))\;\;\mbox{for all}\;\;\mathbf x, \mathbf y \in \F^n, \lambda \in \F.
\end{multline}
Since Lemma~\ref{lem:condimpliesphi} implies that
\begin{equation}\label{thm:homogenphipsi:eq8}
P(\mathbf x) = P(\phi(\mathbf x))\;\;\mbox{for all}\;\;\mathbf x \in \F^n,
\end{equation}
by subtracting~\eqref{thm:homogenphipsi:eq8} from~\eqref{thm:homogenphipsi:eq7} we obtain that
\begin{equation*}
\lambda P(\mathbf y) = \lambda P(\psi(\mathbf y))\;\;\mbox{for all}\;\;\mathbf y \in \F^n, \lambda \in \F.
\end{equation*}
In particular, for $\lambda = 1$ we obtain that
\begin{equation*}
P(\mathbf y) = P(\psi(\mathbf y))\;\;\mbox{for all}\;\;\mathbf y \in \F^n.
\end{equation*}
Thus, $\psi$ preserves $P$.

\paragraph{Case (III): $\deg(P) > 1$} Then Claim~\ref{thm:homogenphipsi:part1} implies that 
\begin{equation}\label{thm:homogenphipsi:eq9}
\phi(\mathbf 0) \in \rad(P).
\end{equation}
Since $\phi$ and $\psi$ satisfy the condition~\eqref{thm:homogenphipsi:eqq1},
\begin{multline}\label{thm:homogenphipsi:eq10}
 P(\phi(\mathbf 0) + \psi(\mathbf x)) = P(\phi(\mathbf 0) + 1\cdot \psi(\mathbf x)) = P(\mathbf 0 + 1\cdot \mathbf x)\\ = P(\mathbf 0 + \mathbf x) = P(\mathbf x)\;\;\mbox{for all}\;\;\mathbf x \in \F^n.
\end{multline}
Therefore,
\[
P(\psi(\mathbf x)) \overset{\eqref{thm:homogenphipsi:eq9}}{=\joinrel=}  P(\phi(\mathbf 0) + \psi(\mathbf x))\overset{\eqref{thm:homogenphipsi:eq10}}{=\joinrel=} P(\mathbf x)\;\;\mbox{for all}\;\;\mathbf x \in \F^n,
\]
which means that $\psi$ preserves $P$.
\paragraph{Claim~\ref{thm:homogenphipsi:part3}}%. $\psi_{\rad(P)}$ preserves $P_{\rad(P)}$}
The proposition that $\psi$ preserves $P$ could be rewritten as follows
\begin{equation}\label{thm:homogenphipsi:eq1001}
P \circ \psi = P.
\end{equation}
In addition,
\begin{equation}\label{thm:homogenphipsi:eq11}
P_{\rad(P)} \circ  \pi_{\rad(P)} = P
\end{equation}
by the definition of $P_{\rad(P)}$ and 
\begin{equation}\label{thm:homogenphipsi:eq12}
\psi_{\rad(P)} \circ \pi_{\rad(P)} = \pi_{\rad(P)} \circ \psi
\end{equation}
by the definition of $\psi_{\rad(P)}$. Hence,
\begin{multline}\label{thm:homogenphipsi:eq13}
\left(P_{\rad(P)} \circ  \psi_{\rad(P)}\right) \circ \pi_{\rad(P)} = P_{\rad(P)} \circ  \left(\psi_{\rad(P)}\circ \pi_{\rad(P)}\right)\\
\overset{\eqref{thm:homogenphipsi:eq12}}{=\joinrel=} P_{\rad(P)} \circ  \left(\pi_{\rad(P)} \circ \psi \right)
= \left(P_{\rad(P)} \circ  \pi_{\rad(P)} \right) \circ \psi\\ \overset{\eqref{thm:homogenphipsi:eq11}}{=\joinrel=} P \circ \psi \overset{\eqref{thm:homogenphipsi:eq1001}}{=\joinrel=} P \overset{\eqref{thm:homogenphipsi:eq11}}{=\joinrel=} P_{\rad(P)} \circ  \pi_{\rad(P)}.
\end{multline}
Since $\pi_{\rad(P)}$ is a surjective map, then~\eqref{thm:homogenphipsi:eq13} implies that
\[
P_{\rad(P)} \circ  \psi_{\rad(P)} = P_{\rad(P)}.
\]
Thus, $\psi_{\rad(P)}$ preserves $P_{\rad(P)}$.


\paragraph{Claim~\ref{thm:homogenphipsi:part4}}%. $\pi_{\rad(P)} \circ \phi = \pi_{\rad(P)} \circ \psi$}
This is equivalent to the following
\begin{equation*}\label{thm:homogenphipsi:eq14}
\left(\phi(\mathbf x) - \psi(\mathbf x)\right) \in \rad(P)\;\;\mbox{for all}\;\;\mathbf x \in \F^n,
\end{equation*}
which is equivalent to that
\begin{equation}\label{thm:homogenphipsi:eq15}
P(\mathbf a + \lambda \left(\phi(\mathbf x) - \psi(\mathbf x)\right)) = P(\mathbf a)\;\;\mbox{for all}\;\;\mathbf a, \mathbf x \in \F^n, \lambda \in \F.
\end{equation}
Let $\mathbf a, \mathbf x \in \F^n$ and $\lambda  \in \F.$ Then
{\setlength{\jot}{15pt}
\begin{multline*}\label{thm:homogenphipsi:eq16}
P\left(\mathbf a + \lambda\left(\phi(\mathbf x) - \psi(\mathbf x)\right)\right)\\
\overset{\scriptsize\mbox{Definition of $P_{\rad(P)}$}}{=\joinrel=\joinrel=\joinrel=\joinrel=\joinrel=\joinrel=\joinrel=\joinrel=\joinrel=\joinrel=\joinrel=}
P_{\rad(P)} \circ \pi_{\rad(P)}\left[\mathbf a + \lambda\left(\phi(\mathbf x) - \psi(\mathbf x)\right)\right]\\
\overset{\scriptsize\mbox{$\pi_{\rad(P)}$ is linear}}{=\joinrel=\joinrel=\joinrel=\joinrel=\joinrel=\joinrel=\joinrel=\joinrel=\joinrel=\joinrel=\joinrel=}
P_{\rad(P)} \left[\pi_{\rad(P)}(\mathbf a) + \pi_{\rad(P)}( \lambda \phi(\mathbf x)) - \lambda\pi_{\rad(P)}( \psi(\mathbf x))\right]\\
\overset{\scriptsize\mbox{$\psi_{\rad(P)}$ is invertible}}{=\joinrel=\joinrel=\joinrel=\joinrel=\joinrel=\joinrel=\joinrel=\joinrel=\joinrel=\joinrel=\joinrel=}
P_{\rad(P)} \left[\psi_{\rad(P)}\left(\psi_{\rad(P)}^{-1}\left(\pi_{\rad(P)}(\mathbf a)\right)\right)\right.\phantom{XXXXXXXXX}\\\phantom{XXXXXXXXXXXXXXXXXX}\left.
+ \pi_{\rad(P)}(\lambda  \phi(\mathbf x)) - \lambda \pi_{\rad(P)}( \psi(\mathbf x))\right]\\
\overset{\scriptsize\mbox{Definition of $\psi_{\rad(P)}$}}{=\joinrel=\joinrel=\joinrel=\joinrel=\joinrel=\joinrel=\joinrel=\joinrel=\joinrel=\joinrel=\joinrel=}
P_{\rad(P)} \left[\psi_{\rad(P)}\left(\psi_{\rad(P)}^{-1}\left(\pi_{\rad(P)}(\mathbf a)\right)\right)\right.\phantom{XXXXXXXXX}\\\phantom{XXXXXXXXXXXXXXX}\left. + \pi_{\rad(P)}( \lambda \phi(\mathbf x)) - \lambda \psi_{\rad(P)}\left(\pi_{\rad(P)}(\mathbf x)\right)\right]\\
\overset{\scriptsize\mbox{Rearrangement of summands}}{=\joinrel=\joinrel=\joinrel=\joinrel=\joinrel=\joinrel=\joinrel=\joinrel=\joinrel=\joinrel=\joinrel=\joinrel=\joinrel=\joinrel=\joinrel=\joinrel=}
P_{\rad(P)} \left[\psi_{\rad(P)}\left(\psi_{\rad(P)}^{-1}\left(\pi_{\rad(P)}(\mathbf a)\right)\right)\right.\phantom{XXXXXXX}\\\phantom{XXXXXXXXXXXXXXX}\left. - \lambda \psi_{\rad(P)}\left(\pi_{\rad(P)}(\mathbf x)\right) + \pi_{\rad(P)}( \lambda \phi(\mathbf x)) \right]\\
\overset{\scriptsize\mbox{$\psi_{\rad(P)}$ is linear}}{=\joinrel=\joinrel=\joinrel=\joinrel=\joinrel=\joinrel=\joinrel=\joinrel=\joinrel=\joinrel=\joinrel=}
P_{\rad(P)} \left[\psi_{\rad(P)}\left(\psi_{\rad(P)}^{-1}\left(\pi_{\rad(P)}(\mathbf a)\right) - \lambda \pi_{\rad(P)}(\mathbf x)\right)\right.\phantom{XXX}\\\phantom{XXXXXXXXXXXXXXXXXXXXXXXXXX}\left. + \pi_{\rad(P)}( \lambda \phi(\mathbf x)) \right]\\
\overset{\scriptsize\mbox{Lemma~\ref{lem:TwoCondLift}}}{=\joinrel=\joinrel=\joinrel=\joinrel=\joinrel=\joinrel=\joinrel=\joinrel=\joinrel=\joinrel=\joinrel=}
P_{\rad(P)} \left[\left(\psi_{\rad(P)}^{-1}\left(\pi_{\rad(P)}(\mathbf a)\right) - \lambda \pi_{\rad(P)}(\mathbf x)\right) + \pi_{\rad(P)}( \lambda \mathbf x) \right]\\
\overset{\scriptsize\mbox{$\pi_{\rad(P)}$ is linear}}{=\joinrel=\joinrel=\joinrel=\joinrel=\joinrel=\joinrel=\joinrel=\joinrel=\joinrel=\joinrel=\joinrel=}
P_{\rad(P)} \left[\psi_{\rad(P)}^{-1}\left(\pi_{\rad(P)}(\mathbf a)\right) - \lambda \pi_{\rad(P)}(\mathbf x) + \lambda \pi_{\rad(P)}( \mathbf x) \right]\\ 
=P_{\rad(P)} \left[\psi_{\rad(P)}^{-1}\left(\pi_{\rad(P)}(\mathbf a)\right) \right]\phantom{XXXXXX}\\
\overset{\scriptsize\mbox{$\psi_{\rad(P)}$ preserves $P_{\rad(P)}$}}{=\joinrel=\joinrel=\joinrel=\joinrel=\joinrel=\joinrel=\joinrel=\joinrel=\joinrel=\joinrel=\joinrel=\joinrel=\joinrel=\joinrel=}
P_{\rad(P)} \circ \psi_{\rad(P)} \left[\psi_{\rad(P)}^{-1}\left(\pi_{\rad(P)}(\mathbf a)\right) \right]\\
= P_{\rad(P)} \circ \pi_{\rad(P)}(\mathbf a)
\overset{\scriptsize\mbox{Definition of $P_{\rad(P)}$}}{=\joinrel=\joinrel=\joinrel=\joinrel=\joinrel=\joinrel=\joinrel=\joinrel=\joinrel=\joinrel=\joinrel=}
P(\mathbf a).
\end{multline*}}
Thus, \eqref{thm:homogenphipsi:eq15} holds for all $\mathbf a, \mathbf x \in \F^n$, $\lambda  \in \F$ and consequently $\pi_{\rad(P)} \circ \phi = \pi_{\rad(P)} \circ \psi$.
\paragraph{Claim~\ref{thm:homogenphipsi:part5}}%. $\psi_{\rad(P)}$ is a unique linear map satisfying the conditions of the lemma}
Indeed, the map $\psi_{\rad(P)}$ is linear by Lemma~\ref{lem:PPhiPsi}, preserves $P_{\rad(P)}$ by the claim~\ref{thm:homogenphipsi:part3} and the following equalities hold
\[
\pi_{\rad(P)} \circ \phi 
\overset{\scriptsize\mbox{Claim~\ref{thm:homogenphipsi:part4}}}{=\joinrel=\joinrel=\joinrel=\joinrel=\joinrel=}
\pi_{\rad(P)} \circ \psi
\overset{\scriptsize\mbox{Lemma~\ref{lem:PPhiPsi}}}{=\joinrel=\joinrel=\joinrel=\joinrel=\joinrel=\joinrel=}
\psi_{\rad(P)} \circ \pi_{\rad(P)}. 
\]
The last equality implies that $\psi_{\rad(P)}$ is unique by Lemma~\ref{lem:PPhiPsi}.
\end{proof}

\end{document}
